

\begin{obeassessment}
\textbf{Kuis Bab 1}
1. Apa yang dimaksud dengan \textit{non-repudiation} dan mengapa hal ini sangat penting dalam transaksi elektronik?
2. Jika seseorang mengklaim bahwa enkripsi sama dengan hashing, jelaskan mengapa pernyataan tersebut salah berdasarkan tujuan utamanya.
3. Tuliskan kembali sistem kriptografi Caesar cipher sebagai quintuple
   $(\mathcal{E},\mathcal{D},\mathcal{M},\mathcal{K},\mathcal{C})$, lengkap dengan
   himpunan dan fungsinya.
4. Sebuah tim pengembang hendak mengamankan komunikasi antara 200 cabang bank.
   Hitung berapa kunci yang diperlukan bila memakai kriptografi simetri, dan berapa
   pasangan kunci bila memakai kriptografi nir-simetri. Jelaskan mengapa solusi
   praktisnya justru menggabungkan keduanya.
5. Jelaskan perbedaan antara \textit{encryption at rest} dan \textit{encryption in
   motion}, serta sebutkan satu contoh penerapan nyata untuk masing-masing.
\end{obeassessment}
\begin{competencychecklist}
\begin{tabularx}{\linewidth}{|X|c|}
\hline
Indikator Kompetensi & Tercapai \\
\hline
Saya dapat menjelaskan definisi dan terminologi dasar kriptografi & [ ] \\
\hline
Saya dapat membedakan 4 layanan utama kriptografi (CIA + Non-repudiation) & [ ] \\
\hline
Saya dapat membedakan algoritma simetri, nir-simetri, dan fungsi hash & [ ] \\
\hline
Saya dapat menjelaskan Prinsip Kerckhoffs beserta konsekuensinya & [ ] \\
\hline
Saya dapat memberikan contoh kasus urgensi kriptografi di dunia TI & [ ] \\
\hline
\end{tabularx}
\end{competencychecklist}
